\documentclass[11pt]{article}
\usepackage[T1]{fontenc}
\usepackage[utf8]{inputenc}
\usepackage{lmodern}
\usepackage[letterpaper,margin=0.84in,headheight=14pt]{geometry}
\usepackage{amsmath,amssymb,amsthm,mathtools}
\usepackage{booktabs,array,longtable,enumitem,microtype,xcolor}
\usepackage{xurl,hyperref,fancyhdr,listings}
\hypersetup{colorlinks=true,linkcolor=teal!60!black,urlcolor=teal!60!black,citecolor=teal!60!black,
 pdftitle={One proof per mechanism: simplifying the K215 correctness proofs},
 pdfauthor={Research and implementation proposal for ProveIt}}
\setlength{\parskip}{0.42em}
\setlength{\parindent}{0pt}
\setlength{\emergencystretch}{2em}
\setlist{itemsep=0.25em,topsep=0.4em}
\pagestyle{fancy}
\fancyhf{}
\fancyhead[L]{\small ProveIt / K215 proof simplification}
\fancyhead[R]{\small 22 October 2026}
\fancyfoot[C]{\thepage}
\newcommand{\code}[1]{\texttt{\detokenize{#1}}}
\newcommand{\Halts}{\operatorname{Halts}}
\newcommand{\Reach}{\operatorname{Reach}}
\newcommand{\Con}{\operatorname{Con}}
\newcommand{\Imp}{\mathbin{\Rightarrow}}
\newtheorem{theorem}{Theorem}[section]
\newtheorem{proposition}[theorem]{Proposition}
\newtheorem{lemma}[theorem]{Lemma}
\lstset{basicstyle=\small\ttfamily,breaklines=true,columns=fullflexible,
 keepspaces=true,frame=single,rulecolor=\color{black!15},backgroundcolor=\color{black!2},
 xleftmargin=0.5em,xrightmargin=0.5em,aboveskip=0.7em,belowskip=0.7em}
% Final source accounting and accepted combined-check scopes.
\newcommand{\RocqImplicationSummary}{%
The Rocq implementation consumes the exact same six-goal parser and proof
data as the Lean shared-basis producer. It adapts the earlier generic Rocq
proof interface from the repository's shared-proof-core proposal
\cite{sharedproposal}, replacing the maintained concrete certificate
interpreter/replay with a generic theorem and a small numeric adapter.
The generic module has 298 lines / 15,004 bytes, and the adapter has
55 lines / 2,445 bytes. Together they replace 1,169 lines / 22,673 bytes
with 353 lines / 17,449 bytes: 816 lines and 5,224 bytes saved, including
the entire new core. This is a 23.0\% reduction in that selected proof source.

The adapter preserves all three original \code{DomImp} fields and the
\code{rho8_all} contract. Its stronger \code{rho8_all_of_restricted} result
also exposes that the truth field is unused by this particular implication
theorem. The six generic exports, this numeric specialization, and the
existing seed equivalence have nine closed assumption reports in the
fresh combined compatibility build.

This optional integration patch follows the Rocq scheduler patch because
both update the same producer and manifest. It adds a Rocq renderer and
adapter template, updates generated sources and compile order, and depends
on the shared six-goal proof data supplied by the Lean patch. A producer
replay reproduces all five changed Rocq modules byte-for-byte. It exercises
the calculus producer, not full machine regeneration.

With both Rocq refactors, the three old concrete proof modules total
40,773 bytes / 1,579 lines; the five replacement modules total
39,074 bytes / 864 lines. The net saving is 1,699 bytes (4.2\%) and
715 lines. Template copies, generator code, and audit artifacts are outside
this explicitly stated proof-module comparison.}
\newcommand{\CombinedVerificationSummary}{%
A final combined Lean build freshly compiled all 22 supplied modules in one
isolated tree. It uses the restored six-goal basis, its rebuilt numeric
provider, all four scheduler clients, fourteen shared-basis typed clients,
ten scheduler typed clients, and the quantifier counterexample. No compiled
project object was reused. All modules passed; only the six existing
\code{TarskiComplete} warnings remained. This checks the interaction between
the two refactors, without extending the claim to the full machine endpoint.

The final combined Rocq check freshly compiled twelve dependency/component
modules and a typed client, then passed recursive \code{coqchk}. Its nine
audited implication/adapter/seed results were closed under the global context.
The self-contained driver's separate fresh run confirmed the same result using
the included original source bytes and the documented compatibility copies.

Finally, all four production patches were applied in order to an isolated
minimal baseline. The integration check matched thirteen checked proof-source
hashes and reproduced three Lean basis outputs, two enumerator variants and
five Rocq outputs exactly. The actual repository sources remained unchanged.
This validates source correspondence and producer integration as well as the
separate proof checks.}


\title{\vspace{-1.0em}\textbf{One proof per mechanism}\\[0.35em]
\large Shared implication proofs and a checked candidate scheduler\\
\large for the 215-state ZFC-consistency machine}
\author{A source-pinned proposal and implementation packet for ProveIt}
\date{22 October 2026}

\begin{document}
\maketitle
\vspace{-1em}
\begin{abstract}
The immediate recommendation is to share the implication derivation across its
numeric and unrestricted interpretations, then share the candidate-search
argument across enumerators. These are mathematically independent changes: the first provides
a substantial source reduction; the second reduces duplicated reasoning and
weakens the scheduler's assumptions, with little immediate aggregate size gain.
The supplied Lean patch reduces the two implication proof files from 811 to 401
lines. The scheduler patch replaces sixteen separately maintained inductive
proofs in four Lean enumerators by three shared inductive proofs. A constructive
Rocq scheduler and current-source adapters are also supplied.

The packet contains complete source overlays, reviewable patches, reproduction
scripts, source hashes, compiler logs, and a checked counterexample to an unsafe
quantifier weakening. It builds on ideas and derivations already present in the
repository, explicitly distinguishing their reuse from the new integration and
verification work. Component checks are fresh; a new complete physical-machine
endpoint build is not claimed. Machine tables and the existing consistency
statement are outside the proposed changes.
\end{abstract}

\section{Recommendation and exact scope}
\label{sec:recommendation}

\textbf{Adopt the shared implication basis first.} It is the clearest immediate
win: one derivation serves both the dominated numeric calculus and the
unrestricted implicational calculus. The supplied patch preserves the public
interfaces, includes the maintained producer, and has been checked against the
actual retained sources.

\textbf{Adopt the generic scheduler for reuse across candidates.} Its contract
mentions only a candidate index, a decreasing pending-work rank, and two
successful-step laws. This isolates the mathematical search argument from the
encoding and the prover state. For a standalone K215-only source-size target,
retaining the smaller concrete argument is a defensible choice until more
clients share the core or obsolete compatibility wrappers can be retired.

\textbf{Then address decoder certificates.} A route and its internal decoder
clock should be selected before the tape state. This is a promising second
stage, already identified by earlier repository work. It needs its own
implementation and full coverage checks; it is not part of the checked patches
in this packet.

\subsection{Pinned baseline, existing work, and measurement boundary}

All comparisons use repository commit
\begin{center}
\code{57541fe0cfe3af62b4cb58747d8d1d7c16b338d5}.
\end{center}
The actual phase-twelve endpoint imports 179 project modules, comprising
2,129,439 source bytes; the packet independently inventories that closure.
Those are baseline measurements, not the results of the component refactor
\cite{phase12}. In particular, a recent commit message is not a reliable
description of the files retained at this revision.

The six-goal shared implication basis already appears in commits
\code{3a7a61ee} and \code{49e99c5b}. The retained phase-twelve proof files still
use the older three-goal generic core and the separate unrestricted derivation.
This packet reconciles the existing shared basis with that actual baseline,
synchronizes its producer, and freshly checks the relevant clients
\cite{sharedhistory,phase12}. It does not claim to discover that derivation.
Similarly, arbitrary-state completeness, cofinal encodings, and scheduler
factoring have antecedents in the repository's checked-transport and
further-interface proposals \cite{checkedtransport,furtherinterface}.

\begin{table}[ht]
\centering
\small
\begin{tabular}{p{0.47\linewidth}rrr}
\toprule
Selected production proof sources & Before & After & Change\\
\midrule
Lean implication pair, lines & 811 & 401 & $-410$\\
Lean implication pair, bytes & 42,454 & 18,243 & $-24,211$\\
Lean four enumerators + new core, lines & 2,666 & 2,659 & $-7$\\
Lean four enumerators + new core, bytes & 109,663 & 109,339 & $-324$\\
Rocq two scheduler clients + new core, lines & 410 & 511 & $+101$\\
Rocq two scheduler clients + new core, bytes & 18,100 & 21,625 & $+3,525$\\
Rocq implication proof + new core, lines & 1,169 & 353 & $-816$\\
Rocq implication proof + new core, bytes & 22,673 & 17,449 & $-5,224$\\
\bottomrule
\end{tabular}
\caption{Comments and whitespace are included. Each row has its stated source
boundary; these are neither whole-repository sizes nor compiled proof sizes.}
\label{tab:metrics}
\end{table}

The complete six-file Lean implication integration, including proof data,
renderer, producer and new adapter template, changes 80,277 to 55,874 bytes:
24,403 bytes saved. Audit programs, historical evidence, documentation, and the
new standalone reproduction bundles are separate deliverable overhead. The
source ledgers record every before/after hash, so these boundaries are
reviewable rather than inferred from a hand-picked proof fragment.

\subsection{Preserve the precise consistency theorem}

The maintained Lean endpoint states \cite{endpoint}
\[
 \code{BlankHalts machine215}
 \quad\longleftrightarrow\quad
 \code{BProv ZFCax_s [] Form.fBot}.
\]
Its nonhalting equivalence is obtained by negating both sides. Thus the operative
consistency predicate is syntactic non-refutability for the repository's chosen
formalization of ZFC. The objective here is to shorten the proof of that
equivalence while preserving the syntax, seed, program, compiler relation, and
literal machine.

There are two different formalization obligations. The existing external
proof-assistant endpoint establishes the displayed equivalence in its ambient
formal environment. A theorem stated \emph{inside ZFC}, or inside a weaker
arithmetical metatheory, additionally needs the appropriate internal coding and
derivability bridge. The repository itself explicitly leaves an internal PA
endpoint open \cite{overview}. No component refactor in this packet fills that
gap or asserts consistency or unconditional nonhalting. These distinctions
should remain in the repository's theorem descriptions after integration.

\section{Share the implication derivation at its weakest interface}
\label{sec:implication}

The numeric side conditions obscure a simpler fact: the derivation only needs
a relation that is reflexive and remains true when its right-hand object is
embedded in either operand of implication. Arithmetic, parsing, truth values,
and injectivity are unnecessary in this core.

Let $A$ be a type with an operation $\Imp:A\times A\to A$ and a relation
$\preceq$ satisfying
\begin{align*}
 a&\preceq a,\\
 u\preceq x&\Longrightarrow u\preceq(x\Imp y),\\
 u\preceq y&\Longrightarrow u\preceq(x\Imp y).
\end{align*}
For a provability predicate $T:A\to\mathsf{Prop}$, assume modus ponens and the
restricted row
\[
 u\preceq a\quad\Longrightarrow\quad
 T\bigl(((c\Imp u)\Imp a)\Imp((a\Imp c)\Imp(b\Imp c))\bigr).
\]
Write the formula in parentheses as $\rho_8(a,b,c,u)$. A single
condensed-detachment graph supplies the following exports:
\begin{enumerate}
\item unrestricted $\rho_8$: $T(\rho_8(a,b,c,u))$ for all four arguments;
\item composition: $T((a\Imp b)\Imp((b\Imp c)\Imp(a\Imp c)))$;
\item lifting under an antecedent: $T(a\Imp b)\Rightarrow
      T((x\Imp a)\Imp(x\Imp b))$;
\item $K$: $T(a\Imp(b\Imp a))$;
\item Peirce: $T(((a\Imp b)\Imp a)\Imp a)$;
\item $S$: $T((a\Imp(b\Imp c))\Imp((a\Imp b)\Imp(a\Imp c)))$.
\end{enumerate}

The inherited proof-data file describes a 58-step derivation with six goals.
Its Lean rendering has 56 named private lemmas, with the remaining final nodes
inlined into public proofs. Its generator emits ordinary proof terms. The generic theory assumes neither
transitivity of $\preceq$ nor implication injectivity. This is useful beyond
cosmetic generality: it prevents numeric coding lemmas from becoming accidental
dependencies of the propositional completeness bridge.

\subsection{Two interpretations, one maintained proof graph}

For the \textbf{numeric interpretation}, take the encoded implication operation
and $u\preceq a$ to mean $u\leq a$. The existing pairing inequalities discharge
the two growth properties. The concrete restricted calculus supplies modus
ponens and its dominated $\rho_8$ rule. The adapter's unrestricted-row,
composition and lifting theorems are then direct applications of the core. The
Lean numeric adapter is byte-identical to the baseline in the supplied shared
basis patch.

For the \textbf{unrestricted interpretation}, take ordinary implication syntax
and define $u\preceq a$ to be $\mathsf{True}$. All three relation obligations are
trivial, and the unrestricted axiom supplies every restricted row. The existing
\code{Lukasiewicz} namespace can keep its named results while deriving $K$, $S$,
Peirce and composition from the same core. This changes
\code{Lukasiewicz.lean} from 524 lines / 28,731 bytes to 101 lines / 4,097 bytes.
The generic core grows modestly, from 287 to 300 lines, because its public
interface now includes the three additional basis results.

\textbf{Falsity remains a separate ingredient.} The numeric and Tarski bridges
still need their original falsity or $\rho_6$ assumptions. Pure implication
alone is not being claimed to provide ex falso for an unrelated constant.
Likewise, the existing \code{DomImp} interface should retain its truth-related
field for compatibility unless a separate consumer census justifies changing
it. Its presence in the adapter does not force the generic derivation to use it.

\subsection{Producer consistency is part of the refactor}

Copying shorter Lean files without updating the producer would leave the
repository in an unstable state: the next regeneration would restore the old
proofs or fail a pinned-hash check. The supplied patch therefore changes six
production paths together: the two Lean proof files, the maintained producer,
the condensed-detachment renderer and six-goal proof data, and the small
unrestricted adapter template. The producer preserves its existing
\code{--out}, \code{--check}, and default command-line behavior.

Fresh producer replay reproduces all three maintained Lean outputs exactly.
The four legacy renderer modes also reproduce their previous bytes. In
particular, adding the shared Lean interpretation does not silently replace a
different legacy Rocq rendering. The generator and its hash checks are useful
engineering tools, but correctness comes from checking the emitted proofs in
the assistant. No trust in the search program or renderer is added.

The Lean audit has fourteen explicitly typed old-interface clients. The six
generic exports and unrestricted $K$, $S$, Peirce and composition are axiom-free.
Numeric consumers retain \code{propext} and \code{Quot.sound}; the Tarski
completeness consumers retain the standard classical profile. No new axiom or
\code{sorryAx} appears in the component evidence.

\subsection{Rocq implication component}

\RocqImplicationSummary

\section{Prove candidate scheduling once}
\label{sec:scheduler}

The next duplicated mechanism is finite progress through an infinite sequence
of finite candidates. Its proof should not know the formula representation,
the stack discipline, or the compiler's command constructors. It needs only a
small progress interface and a separate adequacy lemma for a selected target.

\subsection{Operational interface and progress laws}

Let $S$ be a state type and
$\mathsf{step}:S\to\mathsf{Option}(S)$ a partial deterministic transition.
Retain the existing concrete run function with equations
\begin{align*}
 \mathsf{run}(0,s)&=\mathsf{some}(s),\\
 \mathsf{run}(n+1,s)&=
 \begin{cases}
   \mathsf{run}(n,t),&\mathsf{step}(s)=\mathsf{some}(t),\\
   \mathsf{none},&\mathsf{step}(s)=\mathsf{none}.
 \end{cases}
\end{align*}
Define halting by a finite successful prefix followed by an undefined step:
\[
 \Halts(s)\;:\!\iff\;
 \exists n,t,\quad \mathsf{run}(n,s)=\mathsf{some}(t)
                  \;\land\;\mathsf{step}(t)=\mathsf{none}.
\]
Let $p,i:S\to\mathbb N$ be a pending-work rank and the candidate index. A
boundary is a state with $p(s)=0$. Require only these laws for successful steps
$s\to t$:
\begin{align}
 p(s)=0 &\Longrightarrow i(t)=i(s)+1,\label{eq:advance}\\
 p(s)\ne0 &\Longrightarrow p(t)<p(s)\ \land\ i(t)=i(s).
 \label{eq:descend}
\end{align}

There is \emph{no bound on the rank loaded by a boundary step}. In particular,
the proof does not require $p(t)\leq i(t)$, monotone encodings, candidate
injectivity, or an explicit time budget. The rank is a progress measure, not
necessarily a literal instruction count.

\begin{theorem}[Reach the target boundary or halt first]
\label{thm:seek}
For any $s$ and $K$ with $i(s)\leq K$, either $\Halts(s)$ or there exist $n,t$
such that
\[
 \mathsf{run}(n,s)=\mathsf{some}(t),\qquad p(t)=0,\qquad i(t)=K.
\]
\end{theorem}

\begin{proof}
Use lexicographic induction on $(K-i(s),p(s))$, retaining $i(s)\leq K$.
If $p(s)=0$ and $i(s)=K$, use the zero-step run. Otherwise inspect the step.
An undefined step proves halting immediately. At a successful active step,
\eqref{eq:descend} preserves the first component and strictly decreases the
second. At a successful boundary step below the target,
\eqref{eq:advance} strictly decreases the first component and preserves the
bound $i(t)\leq K$; the new second component may be arbitrarily large.
Apply the induction hypothesis and prepend the step to its finite witness.
\end{proof}

Lean implements this as \code{reach_or_halt} with an explicit lexicographic
termination measure. Rocq implements \code{seek} by nested strong induction
using an equivalent nonnegative offset. The argument is constructive; the
Rocq component's assumption audit is closed under the global context.

\subsection{Use a local completion theorem with the correct quantifiers}

\begin{theorem}[Completion at a fixed reachable boundary]
\label{thm:complete}
Suppose $i(s)\leq K$ and
\[
 \forall n,t,\quad
 \bigl(\mathsf{run}(n,s)=\mathsf{some}(t)\land p(t)=0\land i(t)=K\bigr)
 \ \Longrightarrow\ \Halts(t).
\]
Then $\Halts(s)$.
\end{theorem}

\begin{proof}
Apply Theorem~\ref{thm:seek}. Its halting alternative already gives the result.
Otherwise the hypothesis gives halting from the reached boundary. Concatenate
the successful prefix with that boundary's halting witness.
\end{proof}

This is the supplied \code{complete_at} interface. It requests adequacy only
for boundary payloads actually reachable from the given start, which can be
strictly weaker than adequacy for all states. The target $K$ is selected before
the universal quantification over those payloads. One useful stronger
corollary assumes cofinally many indices that work for every boundary payload:
\[
 \forall N\;\exists K\geq N\;\forall t,\quad
   p(t)=0\land i(t)=K\Longrightarrow\Halts(t).
\]
That condition yields halting from every state. It is more than the local
interface requires and should be kept as a convenience corollary, not built
into the kernel's record.

\subsection{Instantiate the actual K215 enumerator}

For \code{EnumPerm}, the pending proof-list code and next-candidate index give
the two ranks. Existing concrete transition facts establish
\eqref{eq:advance}--\eqref{eq:descend}. The concrete transition and run
definitions remain unchanged. Lean packages their equations in a
\code{Dynamics} record; Rocq's generic iterator is definitionally equal to the
maintained iterator. Neither needs a second semantic simulation proof.

Given a refutation in \code{C222}, the existing dominated-calculus compiler
produces a command word $P$. For an arbitrary initial state $s$, choose
$Y=i(s)$ and the one candidate
\[
 K=\code{encL}\bigl(\code{ax }0\ 1\ Y,\,
              P\mathbin{+\!+}[\code{drop},\code{nip},\code{nip}]\bigr).
\]
The pairing argument \code{encL_ge} proves $Y\leq K$.
The reset prefix and \code{program_none} make the selected complete word halt
for every register payload supplied at that boundary. Therefore this single
$K$ discharges \code{complete_at}. No search for a candidate adapted to an
unknown future payload is needed. The prefix clears the working formula and
stack fields needed by the compiled contract; it need not reset every register
to a fixed constant. The word's contract already accepts the remaining
arbitrary payload \cite{enumperm}.

The new Lean exports are, schematically,
\begin{lstlisting}
complete_any : C222 F -> HaltsFrom s
halts_iff_any : Inv (toW s) -> (HaltsFrom s <-> C222 F)
\end{lstlisting}
The old boundary-restricted exports remain as compatibility wrappers.
Rocq adds \code{enum_complete_from} and \code{enum_from_iff}, preserving
the old enumerator and seed endpoints.

\textbf{Keep the soundness invariant.} Completeness can start from an arbitrary
payload because the compiled witness resets what it needs. Soundness still
requires the initial formula/cache values to be legitimate. An arbitrary
malformed state might already trigger a halt without representing a ZFC
refutation. Accordingly, the generalized equivalence retains \code{Inv}; the
generalized completeness direction does not require it. Earlier repository
work, including \code{EnumSafeStart254}, already demonstrates related
arbitrary-state completeness \cite{safestart,checkedtransport}.

\section{A checked counterexample protects the interface}
\label{sec:counterexample}

A tempting replacement hypothesis is
\[
 \forall\text{ payload }b\;\forall N\;\exists K\geq N,
 \quad\text{candidate }K\text{ halts when started with payload }b.
\]
This does not suffice if the payload can change between candidates. It puts
the existential target on the wrong side of a universal quantifier.

The packet includes \code{QuantifierCounterexample.lean}, a standalone,
axiom-free proof of the failure. States are pairs $(k,b)$ with
$k\in\mathbb N$ and $b\in\{\mathsf{false},\mathsf{true}\}$. Define
$\pi(0)=\mathsf{false}$ and $\pi(k+1)=\neg\pi(k)$, and set
\[
 \mathsf{step}(k,b)=
 \begin{cases}
 \mathsf{some}(k+1,\neg b),&b=\pi(k),\\
 \mathsf{none},&b\ne\pi(k).
 \end{cases}
\]
Every state is a boundary: $p(k,b)=0$ and $i(k,b)=k$. Thus successful steps
satisfy exactly the scheduler's boundary-advance law; the active-step law is
vacuous.

For each fixed $b$ and lower bound $N$, one of $N,N+1$ has parity different
from $b$, so there is a halting candidate beyond every bound for that fixed
payload. Nevertheless the actual aligned trajectory never halts:
\[
 \mathsf{run}(n,(k,\pi(k)))=
      \mathsf{some}(k+n,\pi(k+n)).
\]
At every candidate the payload has changed to the unique nonhalting parity.
No index kills both possible payloads. The source checks payload-wise
cofinality, aligned nonhalting, and absence of a uniform killing candidate,
with empty axiom profiles for all three main results.

\begin{center}
\begin{tabular}{cccl}
\toprule
Index & Actual payload & Immediate halting payload & Actual next state\\
\midrule
0 & false & true & $(1,\mathsf{true})$\\
1 & true & false & $(2,\mathsf{false})$\\
2 & false & true & $(3,\mathsf{true})$\\
3 & true & false & $(4,\mathsf{false})$\\
\bottomrule
\end{tabular}
\end{center}

This is not a criticism of the maintained K215 compiler: its reset word
provides the stronger, correct quantifier order. The counterexample explains
why that reset/adequacy fact must remain explicit when the scheduler is made
generic. It also prevents a future abstraction from replacing a uniform
target with superficially similar pointwise cofinality.

\section{Integration choices and their actual savings}
\label{sec:integration}

\subsection{Lean: share the induction and keep the concrete interface}

The patch instantiates the scheduler in \code{EnumPerm}, \code{Enum254},
\code{Enum222}, and \code{EnumDom}. Each retains its familiar
\code{run_inv}, \code{run_add}, \code{drain}, and \code{reach} interfaces.
Their proofs invoke the shared kernel. The new kernel has three inductive
proofs: run composition, preservation of an invariant along a run, and
reach-or-halt. Halting transport and completion are short consequences.
This replaces sixteen separately maintained inductive arguments across the
four modules by three.

The static source audit preserves 117 original executable declarations and
237 original public theorem statements, comparing declaration text after
comment/whitespace normalization. This is evidence about that explicitly
defined comparison, not a universal semantic-equivalence checker. Actual
proof compilation and typed clients provide the complementary checks.
The existing \code{generate_enum222.py} and \code{generate_enumdom.py} remain
unchanged and regenerate the delivered variants byte-for-byte from the
refactored \code{Enum254}.

The resulting 113-line shared kernel offsets almost all of the removed client
text. Across all four clients, the net reduction is seven lines / 324 bytes.
For \code{EnumPerm} alone plus the new core, the selected source group grows
by 95 lines / 3,411 bytes. The value is a single place to maintain fairness,
fewer proof copies, and a stronger reusable interface. It would be misleading
to describe this step as a large K215 source-size reduction.

\subsection{Rocq: preserve compatibility and synchronize regeneration}

The scheduler component includes \code{CandidateScheduler.v}, a revised
\code{L1_CompPerm.v}, and a shorter \code{SeedPerm.v}. The old
\code{enum_complete}, \code{seedPerm_complete}, \code{seedPerm_sound}, and
\code{seedPerm_iff_thm} types are preserved; generalized arbitrary-state
results are added. Compatibility helpers remain available to downstream
proofs.

The patch updates the generated sources, their \code{rocqperm} templates,
\code{regenerate215.py}, and the extension manifest's hashes and compile
order. Source-only changes without these producer updates would be incomplete.
The two old clients shrink from 18,100 to 17,009 bytes, but the 4,616-byte core
raises the selected total to 21,625 bytes. For a single lane this is an
increase, and the small concrete rank proof remains a reasonable alternative.

There are two routes to later savings: reuse the kernel in additional
enumerators, or retire now-redundant helper wrappers after a whole-consumer
audit. That audit must include generated sources and templates. Preserving
compatibility in this proposal makes its effect easier to review; it is not a
claim that every wrapper deserves permanent retention.

\subsection{Do not mix unrelated size measurements}

Source length, endpoint import closure, elaboration time, proof-object size,
and trusted assumptions measure different things. The packet reports source
counts and scoped compilation evidence. It does not infer faster endpoint
builds from shorter source or smaller proof objects from fewer lines.

In particular, the baseline's 179-module closure is already the result of
earlier dependency factoring. The unrestricted \code{Lukasiewicz} shrinkage
is valuable within that closure, but importing a new scheduler can add a
module even as it removes duplicated induction. Recalculate the closure
after adoption rather than subtracting the four-client aggregate savings
from the K215 endpoint's byte count.

\section{Validation, reproduction, and adoption}
\label{sec:validation}

\subsection{What was actually checked}

The Lean component checks use Lean 4.32.0, compiler commit
\begin{center}\code{8c9756b28d64dab099da31a4c09229a9e6a2ef35}.\end{center}
They need only the bundled
\code{Init}/\code{Std} libraries and the included project providers. They do
not require Mathlib. The shared-basis check freshly compiled fourteen
provider/project modules plus a typed audit. The scheduler check freshly
compiled its six new/refactored/audit modules over fourteen freshly checked
providers. The changed proof files and audit clients emitted no warnings;
the unchanged \code{TarskiComplete} provider emitted six existing linter
warnings, retained in the logs.

The scheduler's generic \code{run_invariant} is axiom-free. Lean's generic
reach and completion results report \code{propext} and \code{Quot.sound};
the concrete audited endpoints use only the existing
\code{propext}, \code{Classical.choice}, and \code{Quot.sound} profile. The
quantifier counterexample's three audited results are axiom-free. These
profiles are measured results of these builds, not a claim that every Lean
encoding of the mathematical argument needs those axioms.

The Rocq scheduler check freshly compiled eleven dependency/component
modules and a typed client, then passed a recursive \code{coqchk}. Six
assumption reports were closed under the global context. The compiler
available for this experiment was Coq 8.18.0 / OCaml 4.14.1. Build copies
replace \code{From Stdlib} with \code{From Coq}; original source bytes and
both hashes are retained. The current production lane specifies Rocq 9.1.0
with rocq-stdlib 9.0.0. Consequently this is a checked compatibility build,
not a fresh claim about that designated toolchain or the whole physical
machine endpoint.

The Rocq providers are the maintained numeric calculus and enumerator
sources. The two older K230 arithmetic files are retained in the
repository's checked-transport packet and match their production origins.
They are not the stub providers from a generic-core demonstration. The
source receipts make that distinction independently inspectable.

\CombinedVerificationSummary

\subsection{How to review the archive}

The top-level \code{README.txt} lists exact commands and patch order.
\code{components/} contains the focused implementation packets, with their
source overlays, generators or adapters, audits, source measurements, and
accepted verification logs. \code{lean/} contains a standalone combined Lean
source bundle and a fresh-build driver; \code{rocq/} provides the analogous
standalone Rocq bundle. \code{integration/} contains the combined patch and
producer check. \code{evidence/} contains the pinned
baseline endpoint inventory and integration checks. The archive manifest
binds the delivered files by SHA-256. Compiled proof objects and failed
development experiments are excluded.

The portable scripts accept explicit compiler and repository paths. They
create isolated output directories and do not require writing into a ProveIt
checkout. Logs preserve the actual commands and execution paths used in
this experiment; portable instructions are separate from those historical
records. Rebuilding the PDF needs an ordinary LaTeX installation with the
packages declared at the top of this source.

\subsection{Suggested adoption sequence}

\begin{enumerate}
\item Start from the pinned revision, or rebase the small patches and repeat
      their source/producer comparisons against the desired target.
\item Adopt the shared implication basis and its producer together. Check the
      old public theorem types and regenerate all maintained outputs.
\item Decide whether the scheduler's cross-candidate reuse justifies its small
      immediate size cost for the intended branch. If so, apply its concrete
      adapters and producer/template changes together.
\item Build the combined affected components from fresh objects. Check their
      axiom profiles and reject any newly admitted lemma or unverified
      provider shortcut.
\item Run the repository's designated Lean and Rocq acceptance gates for
      affected endpoint descendants and both old endpoint clients. Record
      new source/object bindings; retain old checkpoint evidence as history.
\item Only after a consumer census, remove unused compatibility helpers or
      old generated derivations. Measure the actual endpoint closure and
      build characteristics at that point.
\end{enumerate}

No completed source patch should be represented as already certified by an
older whole-endpoint receipt. Conversely, reproducing every expensive machine
site proof is not necessary to assess a clearly labeled component proposal.
The appropriate integration gate supplies that final assurance at adoption.

\section{The next larger simplification: route before tape}
\label{sec:decoder}

After the two checked refactors, the next useful boundary is the decoder
certificate. The current proof already has extensive site coverage; the
phase-twelve checkpoint retains evidence for 6,176 guarded outcomes
\cite{phase12}. The earlier further-interface report proposes choosing a
finite route and a clock before the tape semantics \cite{furtherinterface}.
The opportunity is to package a finite checked decoder path once and reuse a
generic execution theorem, reducing repeated operational reasoning at leaves.

A suitable target statement has the quantifier order
\[
 \forall\text{ source path }p,\quad
 \exists\text{ route }r,\ k,\quad
 \forall\text{ admissible tapes }\tau,\quad
 \mathsf{run}(k,\mathsf{entry}(p,\tau))
       =\mathsf{primitiveEntry}(r,\tau).
\]
This is a design sketch, not a theorem supplied by this packet. The exact
state relation and tape transformation must come from the maintained
semantics. In particular, a proposition of the form
$\forall\tau\,\exists k\,P(\tau,k)$ does not itself justify extracting one
tape-independent $k$. Construct the route and its length from decidable
source-path data, then prove the uniform execution theorem.

Keep the full routed-entry witness and terminal source bit. A physical state
or a short bit suffix alone need not identify the semantic role. The clock
should stop at primitive entry; primitive execution and the return scan have
their own contracts and costs. This prevents an internal decoder bound from
being mistaken for a whole-command running-time bound.

Primitive/return suffix factoring remains open at the pinned checkpoint.
It should be a separate proposal with all guarded outcomes rerun, including
the exception and exit cases. There is no justification for deleting site
obligations merely because many outcomes look similar. For auxiliary lanes
with repeated ad hoc halting budgets, the repository's reachable-suffix
induction proposal is another promising reusable interface, provided it
respects the actual register-equivalence relation rather than adding function
extensionality \cite{furtherinterface}.

The staged plan therefore keeps each review focused: first eliminate a
duplicated implication derivation; then consolidate search progress; then
factor finite decoder execution certificates. The supplied code makes the
first two mechanisms concrete and gives the third a precise proof obligation.

\clearpage
\appendix
\section{Minimal interfaces for maintainers}

The actual Lean scheduler record retains the client's executable definitions:
\begin{lstlisting}
structure Dynamics (S : Type u) where
  step : S -> Option S
  run  : Nat -> S -> Option S
  zero : forall s, run 0 s = some s
  succ : forall n s, run (n + 1) s =
    match step s with
    | some t => run n t
    | none   => none
\end{lstlisting}
Here ASCII arrows are used only for presentation. See the included
\code{CofinalScheduler.lean} for the directly compilable source. The scheduler
laws are separate from the dynamics so ordinary run composition and invariant
transport do not unnecessarily depend on any rank.

The completion API, again in schematic notation, is:
\begin{lstlisting}
complete_at (s : S) (target : Nat)
  (bound : index s <= target)
  (hit : forall n t,
    run n s = some t ->
    pending t = 0 -> index t = target -> Halts t)
  : Halts s
\end{lstlisting}
This is the principal client interface. Exposing an auxiliary proof of the
lexicographic ranking as a requirement on each client would reintroduce the
duplication the abstraction is intended to remove.

\section{Source references}

Repository references below are pinned to the reviewed revision except for the
explicit historical commit. The generated validation receipts in the archive
are the primary evidence for the new component checks and source measurements.

\begin{thebibliography}{99}
\small
\interlinepenalty=10000
\bibitem{endpoint} ProveIt, \emph{KEndpoint215.lean}, exact halting and
nonhalting equivalences.\newline
\url{https://github.com/VladimirReshetnikov/ProveIt/blob/57541fe0cfe3af62b4cb58747d8d1d7c16b338d5/Computability/BusyBeaver/Lean/BusyBeaver/ZFC/KEndpoint215.lean}

\bibitem{phase12} ProveIt, \emph{Twelfth K215 proof-simplification checkpoint},
source closure, checking scope, and remaining work.\newline
\url{https://github.com/VladimirReshetnikov/ProveIt/blob/57541fe0cfe3af62b4cb58747d8d1d7c16b338d5/Computability/BusyBeaver/Research/ZFC/k225-formal/verification/k215/simplification/phase12/README.md}

\bibitem{overview} ProveIt, \emph{K215 proof-simplification overview},
including the external/internal endpoint distinction.\newline
\url{https://github.com/VladimirReshetnikov/ProveIt/blob/57541fe0cfe3af62b4cb58747d8d1d7c16b338d5/Computability/BusyBeaver/Research/ZFC/k225-formal/verification/k215/simplification/README.md}

\bibitem{sharedhistory} ProveIt, historical shared-basis and producer work,
commit \code{49e99c5b7cc6c2e164331961de28bad3904361fd}.\newline
\url{https://github.com/VladimirReshetnikov/ProveIt/commit/49e99c5b7cc6c2e164331961de28bad3904361fd}

\bibitem{sharedproposal} ProveIt, \emph{Shared proof-core proposal}, earlier
representation-independent implication and adapter sources.\newline
\url{https://github.com/VladimirReshetnikov/ProveIt/tree/57541fe0cfe3af62b4cb58747d8d1d7c16b338d5/Computability/BusyBeaver/Research/ZFC/k225-formal/verification/k215/simplification/proposals/shared-proof-core}

\bibitem{checkedtransport} ProveIt, \emph{Checked transport / Rocq seed-rank
packet}, earlier concrete rank and arbitrary-state completeness work.\newline
\url{https://github.com/VladimirReshetnikov/ProveIt/tree/57541fe0cfe3af62b4cb58747d8d1d7c16b338d5/Computability/BusyBeaver/Research/ZFC/k225-formal/verification/k215/simplification/proposals/checked-transport}

\bibitem{furtherinterface} ProveIt, \emph{Further interface proposal},
scheduler extraction, decoder route/clock separation, and suffix induction.\newline
\url{https://github.com/VladimirReshetnikov/ProveIt/blob/57541fe0cfe3af62b4cb58747d8d1d7c16b338d5/Computability/BusyBeaver/Research/ZFC/k225-formal/verification/k215/simplification/proposals/further-interface/report.tex}

\bibitem{enumperm} ProveIt, \emph{EnumPerm.lean}, maintained enumerator,
compiler, and completion proof at the baseline.\newline
\url{https://github.com/VladimirReshetnikov/ProveIt/blob/57541fe0cfe3af62b4cb58747d8d1d7c16b338d5/Computability/BusyBeaver/Lean/BusyBeaver/ZFC/EnumPerm.lean}

\bibitem{safestart} ProveIt, \emph{EnumSafeStart254.lean}, earlier
arbitrary-state completeness interface.\newline
\url{https://github.com/VladimirReshetnikov/ProveIt/blob/57541fe0cfe3af62b4cb58747d8d1d7c16b338d5/Computability/BusyBeaver/Lean/BusyBeaver/ZFC/EnumSafeStart254.lean}

\bibitem{rocqseed} ProveIt, \emph{rocqperm/SeedPerm.v}, retained source for
the seed-completeness adapter.\newline
\url{https://github.com/VladimirReshetnikov/ProveIt/blob/57541fe0cfe3af62b4cb58747d8d1d7c16b338d5/Computability/BusyBeaver/Research/ZFC/k225-formal/rocqperm/SeedPerm.v}
\end{thebibliography}

\end{document}
